
[[1]]
Sei $S^2:=\left\{(x, y, z) \in \mathbb{R}^3: x^2+y^2+z^2=1\right\} \subseteq \mathbb{R}^3$ die Einheitssphäre im $\mathbb{R}^3$. Weiterhin bezeichnen wir den Nordpol mit $N:=(0,0,1)$ und den Südpol mit $S:=(0,0,-1)$. Zu diesen definieren wir die stereographischen Projektionen
$$
\phi_N: S^2 \backslash N \rightarrow \mathbb{R}^2 \quad \text { und } \quad \phi_S: S^2 \backslash S \rightarrow \mathbb{R}^2
$$
durch folgende Vorschrift: Zu $P \in S^2 \backslash N$ gibt es genau eine Gerade $g$ durch $P$ und $N$; wir setzen $(a, b, 0):=\{z=0\} \cap g$ und definieren $\phi_N(P):=(a, b)$. Die Abbildung $\phi_S$ sei analog definiert.

[[1.1]]
Zeigen Sie, dass, nach geeigneten Identifikationen von $\mathbb{R}^2$ mit $\mathbb{C}$, die $\operatorname{Karten}\left(S^2 \backslash N, \phi_N\right)$ und ( $S^2 \backslash S, \phi_S$ ) einen komplexen Atlas von $S^2$ bilden.


[[2]]
In dieser Aufagbe betrachten wir die drei Flächen, die man durch Verkleben von Einheitsquadraten $[0,1]^2 \subset \mathbb{R}^2$ erhält. Dabei werden immer nur parallele Kanten miteinander verklebt. Die Richtung der Markierungen beschreiben die Orientierung des Verklebens.


[[2.1]]
in Quadrat mit vier Kanten. 
Die beiden vertikalen Kanten sind beide nach oben orientiert, 
die beiden horizontalen Kanten beide nach rechts.  
\[
\text{Verklebung: vertikale und horizontale Kanten jeweils parallel und gleich orientiert.}
\]
Das resultierende Flächenobjekt ist topologisch ein \textbf{Torus}.

Zeigen Sie, dass die folgenden Flächen reelle Mannigfaltigkeiten sind, indem Sie einen konkreten Atlas angeben. Definieren diese Atlanten auch eine komplexe Struktur?

[[2.2]]
Ein Quadrat mit vier Kanten. 
Die vertikalen Kanten sind beide nach oben orientiert, 
die obere horizontale Kante nach rechts, 
die untere nach links.  
\[
\text{Verklebung: vertikale Kanten gleich orientiert, horizontale entgegengesetzt orientiert.}
\]
Dies ergibt eine \textbf{Zylinderfläche} 
(bzw.\ ein \textbf{Möbiusband}, falls zusätzlich eine Verdrehung zugelassen wird).

Zeigen Sie, dass die folgenden Flächen reelle Mannigfaltigkeiten sind, indem Sie einen konkreten Atlas angeben. Definieren diese Atlanten auch eine komplexe Struktur?

[[2.3]]
Zwei Quadrate, die eine rechteckige $2\times1$-Fläche bilden.  
Alle vertikalen Kanten sind nach oben orientiert, 
alle horizontalen Kanten nach rechts; 
die gemeinsame Mittelkante zeigt auf beiden Seiten in dieselbe Richtung (nach rechts).  
\[
\text{Verklebung: alle gegenüberliegenden Ränder des gesamten Rechtecks parallel und gleich orientiert.}
\]
Das resultierende Objekt ist topologisch eine \textbf{Fläche vom Geschlecht $2$} (ein \textbf{Doppeltorus}).

Zeigen Sie, dass die folgenden Flächen reelle Mannigfaltigkeiten sind, indem Sie einen konkreten Atlas angeben. Definieren diese Atlanten auch eine komplexe Struktur?
